\documentclass[10pt,a4paper]{amsart}
\usepackage[margin=19mm]{geometry}
\usepackage{amsmath,amssymb,mathtools}
\usepackage[scr=rsfs,cal=euler]{mathalfa}
\usepackage{xfrac}
\usepackage[unicode,hidelinks,bookmarks=false]{hyperref}
\newcommand{\ie}{i.e.\ }
\newcommand{\cf}{cf.\ }
\newcommand{\resp}{respectively\ }
\newcommand{\Subsection}{\subsection}
\providecommand{\nobreakdash}{\nobreak\hbox{-}}
\setlength{\emergencystretch}{2em}
\multlinegap0pt
\usepackage{bm}
\usepackage{bbm}
\usepackage{dsfont}
\usepackage{xspace}
\usepackage{cancel}
\usepackage{mathtools}
\usepackage{cleveref}
\usepackage{amsthm}
\makeatletter
\@ifundefined{theorem}{\newtheorem{theorem}{Theorem}}{}
\@ifundefined{lemma}{\newtheorem{lemma}[theorem]{Lemma}}{}
\makeatother
\DeclareMathOperator{\supp}{\mathrm{supp}}
\newcommand{\A}{\mathcal{A}}
\newcommand{\B}{\mathcal{B}}
\newcommand{\CBP}{\mathrm{CBP}}
\newcommand{\I}{\mathcal{I}}
\newcommand{\R}{\mathbb{R}}
\newcommand{\Z}{\mathbb{Z}}
\renewcommand{\P}{\mathcal{P}}
\title[On the connected blocks polytope]{On the connected blocks polytope}
\author{Justus Bruckamp, Markus Chimani, Martina Juhnke}
\date{Source: arXiv:2304.06318v3 (CC BY 4.0)}
\begin{document}
\maketitle

\noindent\emph{Source and licence.} On the connected blocks polytope. Version arXiv:2304.06318v3 (CC BY 4.0), distributed under the Creative Commons Attribution 4.0 International licence, as recorded in the arXiv licence metadata for this exact version and shown on the abstract page https://arxiv.org/abs/2304.06318v3. Full rights notice: licenses/arxiv-2304.06318-CC-BY-4.0.txt.\par\smallskip

\noindent\textit{Editorial scope.} Counted unit: the Theorem whose source label is \texttt{completefacets} in the article (source0.tex lines 1361--1364, source label \texttt{completefacets}), delivered together with the article's own complete original proof of it (source0.tex lines 1366--1424, one \texttt{proof} environment of 59 lines). No mathematics is written, repaired or re-derived: statement and proof are the article's own text line for line. Required context retained uncounted: combining (lines 1335--1345). Every definition, display and label of those lines is retained unchanged. Omitted from this package and not redistributed: 1--1334, 1346--1360, 1365--1365, 1425--1668. Nothing inside a retained excerpt is omitted or altered: every retained body is delivered exactly as the article's lines are written, so no omission receipt of any kind accompanies this package. Citations: the retained excerpts carry no citation command at all, so no bibliography entry of the article is transcribed and no reference is redistributed. Reference adaptation: none was needed and none was made; every reference inside the delivered unit resolves inside it, so no unresolved reference or citation is produced. Printed slips kept literal: no printed word, symbol, hypothesis or number of the counted statement or of its proof was altered. Layout and class: the article's own class is \texttt{amsart} with packages including hyperref, amssymb, graphicx, color, amsmath, comment, bbm, cleveref, enumerate, mathtools, tikz, soul; the wrapper is the lane's pinned \texttt{amsart} editorial wrapper at 10pt on A4, which restates the theorem-like environments the retained text needs and adds the article's own macro definitions that the retained text needs (\texttt{\textbackslash A}, \texttt{\textbackslash B}, \texttt{\textbackslash CBP}, \texttt{\textbackslash I}, \texttt{\textbackslash P}, \texttt{\textbackslash R}, \texttt{\textbackslash Z}, \texttt{\textbackslash supp}), with extra wrapper packages (bm, bbm, dsfont, xspace, cancel, mathtools, cleveref). The delivered numbering is the unit's own. Assets: no retained span contains a figure environment, a graphics-inclusion command, a PDF-page inclusion, a table or a TikZ picture; no image file is copied into this package, the package declares no asset dependency and no image file is redistributed with it. Licence: version arXiv:2304.06318v3 of this article is distributed under the Creative Commons Attribution 4.0 International licence, as recorded in the arXiv licence metadata for this exact version and shown on the abstract page https://arxiv.org/abs/2304.06318v3. A full rights notice is delivered at licenses/arxiv-2304.06318-CC-BY-4.0.txt. Characters: every retained excerpt is pure ASCII, so the delivered engine, XeLaTeX with its default Latin Modern text font, reports no missing character.
\begin{lemma}
\label{combining}
Let $G$ be a connected graph and $\B$ its set of blocks. Let $(\I^{(1)},\alpha^{(1)}), \ldots, (\I^{(m)},\alpha^{(m)})$ induce independent blocks inequalities for $G$, such that $\B\langle \I^{(i)}\rangle \cap \B\langle \I^{(j)}\rangle= \emptyset$ for all $i,j \in [m]$, $i \neq j$, and $\I \coloneqq \I^{(1)} \cup \ldots \cup \I^{(m)}$ is an independent set of blocks. Moreover, let $\A \subseteq \B$ be a set of blocks, such that $G[\B \setminus \A]$ has exactly $m$ connected components, and each of the connected components contains the blocks $\B\langle \I^{(i)} \rangle$ for exactly one $i \in [m]$. Then there exist blocks $\{A_1,\ldots,A_t\} \subseteq \A$, $t \le m-1$, and weights $\alpha_{A_i} \in \Z_{<0}$, such that $G[\B \setminus \{A_1,\ldots,A_t\}]$ has exactly $m$ connected components $K_1,\ldots,K_m$, with $\B\langle \I^{(i)} \rangle \subseteq K_i$ for all $i \in [m]$, and $(\I,\beta)$, with
\begin{align*}
    \beta_B \coloneqq \begin{cases} \alpha^{(i)}_B, &\text{if } B \in \B\langle \I^{(i)}\rangle, \\
                                    \alpha_{A_i}, &\text{if } B=A_i, \\
                                    0, &\text{otherwise},
    \end{cases}
\end{align*}
defines an independent blocks inequality. Moreover, we have $\sum_{i=1}^t \alpha_{A_i}=-(m-1)$.
\end{lemma}

\begin{theorem}
    \label{completefacets}
    Let $G$ be a connected graph. Then $\CBP(G) = \P(G)$. In particular the independent blocks inequalities, together with nonnegativity constraints, yield the complete and irredundant facet description of $\CBP(G)$.
\end{theorem}

\begin{proof}
    It is clear that $\CBP(G) \subseteq \P(G)$. Let $x \in \P(G)$. We show that there exists a $k \in \mathbb{N}$ and $\lambda_1,\ldots,\lambda_k \in \R_{>0}$, with $\sum_{i=1}^k \lambda_i \le 1$, and vertices $v_1,\ldots,v_k$ of $\CBP(G)$, such that
    \begin{align*}
        x=\sum_{i=1}^k \lambda_iv_i.
    \end{align*}
    This yields that $x$ can be written as a convex combination of the vertices of $\CBP(G)$ since the zero vector can be added to the sum with coefficient $1-\sum_{i=1}^k \lambda_i$. We proceed by induction on the number of blocks in $G$ and show the stronger statement that for every $x \in \P(G)$ there exist $k \in \mathbb{N}$, $\lambda_1,\ldots,\lambda_k \in \R_{>0}$, vertices $v_1,\ldots,v_k$ of $\CBP(G)$, and an independent blocks inequality $(\I,\alpha)$, such that
    \begin{align*}
        x=\sum_{i=1}^k \lambda_iv_i
    \end{align*}
    and
    \begin{align*}
        \sum_{i=1}^k \lambda_i \le \sum_{B \in \B}\alpha_Bx_B.
    \end{align*}
    Since $x \in \P(G)$, this implies $\sum_{i=1}^k \lambda_i \le 1$. Let $\B$ be the set of blocks of $G$. If $\vert \B \vert=1$, the statement is obviously true. Suppose now that $\vert \B \vert=n$ and by induction the statement is true for all connected graphs with less than $n$ blocks. Let $x=(x_{B_1},\ldots,x_{B_n}) \in \P(G)$ and $\mu=\min_{i \in [n]}x_{B_i}$. Then we can write $x$ as
    \begin{align*}
        x=\mu \cdot \mathbf{1}+y,
    \end{align*}
    with $y\coloneqq (x_{B_1}-\mu,\ldots,x_{B_n}-\mu)$. Observe that all entries of $y$ are between $0$ and $1$ and there is at least one entry in $y$ that is $0$. Let $G_y$ be the graph that is given by the union of the blocks of $\supp(y)\coloneqq \{B \in \B \mid y_B > 0\}$ and let $K_1,\ldots,K_m$ be the connected components of $G_y$. For $i\in [m]$ we denote by $y^{(i)}$ the projection of $y$ onto the coordinates $\{B \in \B \mid B \subseteq K_i\}$. We show that $y^{(i)} \in \P(K_i)$ for all $i \in [m]$. Every independent blocks inequality $(\I,\alpha)$ of any $\CBP(K_i)$ induces an independent blocks inequality $(\I,\alpha')$ of $G$, where $\alpha'$ is defined by
    \begin{align*}
        \alpha'_B=\begin{cases}\alpha_B, & \text{if } B \in \B\langle\I\rangle, \\
                                0,& \text{otherwise.}\end{cases}
    \end{align*}
    With this it follows
    \begin{align*}
        \sum_{B \in \B[K_i]}\alpha_By_B^{(i)}&=\sum_{B \in \B}\alpha'_Bx_B-\mu\underbrace{\sum_{B \in \B}\alpha'_B}_{=1} \le 1-\mu\le 1,
    \end{align*}
    which shows that $y^{(i)} \in \P(K_i)$. Hence, by induction, there exist $\lambda_1^{(i)},\ldots,\lambda_{k_i}^{(i)}>0$ and vertices $v_1^{(i)},\ldots,v_{k_i}^{(i)}\in \CBP(K_i)$ and an independent blocks inequality $(\I^{(i)},\alpha^{(i)})$ of $\CBP(K_i)$ such that
    \begin{align*}
        y^{(i)}=\sum_{j=1}^{k_i} \lambda_j^{(i)}v_j^{(i)},
    \end{align*}
    and
    \begin{align}
    \label{lambda}
        \sum_{j=1}^{k_i}\lambda_j^{(i)} \le \sum_{B \in \B[K_i]} \alpha^{(i)}_By^{(i)}_B,
    \end{align}
    Let $w_j^{(i)}\coloneqq v_j^{(i)} \times \{0\}_{B \in \B\setminus\B [K_i]}$. We note that the $w_j^{(i)}$ are vertices of $\CBP(G)$. We can hence write $x$ as
    \begin{align*}
        x=\mu \cdot \mathbf{1} + \sum_{i=1}^m\sum_{j=1}^{k_i} \lambda_j^{(i)}w_j^{(i)}.
    \end{align*}
    For the sum of coordinates we have
    \begin{align*}
        \mu+\sum_{i=1}^m\sum_{j=1}^{k_i} \lambda_j^{(i)} &\overset{\eqref{lambda}}{\le} \mu+\sum_{i=1}^m\sum_{B \in \B[K_i]} \alpha^{(i)}_By_B^{(i)} \\
        &= \mu+\sum_{i=1}^m \left(\sum_{B \in \B[K_i]} \alpha^{(i)}_B(x_B-\mu) \right).
    \end{align*}
    Let $\B(\mu) \coloneqq \{ B \in \B \mid x_B=\mu\}$. Hence, $G[\B \setminus \B(\mu)]=G_y$. By \Cref{combining}, there exists a set of blocks $\A\coloneqq \{A_1,\ldots,A_t\} \subseteq \B(\mu)$, $t \le m-1$, and weights $\alpha_{A_i}$ such that  $(\I,\beta)$, with $\I \coloneqq \bigcup_{i=1}^m\I^{(i)}$, and
    \begin{align*}
        \beta_B\coloneqq \begin{cases} \alpha^{(i)}_B, &\text{if } B \in \B[K_i], \\
                                       \alpha_{A_s}, &\text{if } B=A_s, s \in [t], \\
                                       0, &\text{otherwise}.
        \end{cases}
    \end{align*}
	defines an independent blocks inequality, such that $\sum_{i=1}^t \beta_{A_i}=-(m-1)$. With this we have
    \begin{align*}
        \mu+\sum_{i=1}^m \left(\sum_{B \in \B[K_i]} \alpha^{(i)}_B(x_B-\mu)\right)&=-(m-1)\mu+\sum_{i=1}^m\sum_{B \in \B[K_i]} \alpha^{(i)}_Bx_B \\
        &=\sum_{i=1}^t \beta_{A_i}x_{A_i}+\sum_{i=1}^m\sum_{B \in \B[K_i]} \beta_Bx_B  \\
        &=\sum_{B \in \B}\beta_Bx_B,
    \end{align*}
	which finishes the proof.
\end{proof}

\end{document}
